\section{Complexities and Classes}
Within the recursive class of languages accepted by TMs and hence our
computers, we consider the time and space complexities of their computations.
We so define on an TMs with multiple, but fixed, number of tapes, with finite
tape alphabets. This allows us to divide the class into the \textit{tractable}
and the \textit{intractable} problems.

\subsubsection{Complexity Measures}
In measuring the complexity of a program, we want to find a function $\Phi:
\mathbb{N}\times \mathbb{N}\to \mathbb{N}$ that measures the ``complexity'', or
cost of a TM on some input. We denote $\Phi\left( i,x \right)=\Phi_i\left( x
\right)$ for the complexity of TM $M_i$ with the input $x$. We hence have the
first of Blum's complexity axioms:
\[Dom\left( \phi_i \right)=Dom\left( \Phi_i \right).\]
where $\phi$ is the $i$-th program/TM.  Further, the complexity measure should
be computable for practicality. We have the second of Blum's complexity axioms:
\[\text{The set }\left\{ i,x,B  \in \mathbb{N}^{3}|\Phi_i\left( x \right)
<B\right\}\text{ is recursive}.\]

\begin{definition}[Blum Complexity Measure]
  A \textit{Blum complexity measure} $\iff$ $\Phi\left( i,x \right) $ is a
  partial recursive function that satisfies the two Blum's complexity axioms.
\end{definition}

\subsubsection{Time/Space Complexity Measures}

\begin{definition}[Running Time]
  The \textit{running time} of a TM $M$ on input $\omega$, $Time_M\left(
  \omega \right)$, is the number of steps made by $M$ before halting. If $M$
  does not halt on input $x$, $Time_M\left( x \right) =\infty$.
\end{definition}

\begin{definition}[Time Complexity]
  The \textit{time complexity} of a TM $M$ is the function $T(n)$ that is
  the maximum of the running time of $M$ on  all inputs $\omega$
  of length $n$.
  \begin{center}
    $M$ is $T(n)$ time bounded $\iff\forall$ inputs $\omega$ of length $n$,
    $Time_M\left( \omega \right) \le T(n)$
  \end{center}
\end{definition}

We may further define its analogues for space.
\begin{definition}[Space Consumption]
  The \textit{space consumption} of a TM $M$ on input $\omega$, $Space_M\left(
  \omega \right) $, is the maximum number of cells ``touched'' by $M$ on any of
  its worktapes before halting. If $M$ does not halt on an input $x$,
  $Space_M\left( x \right) =\infty$.
\end{definition}

\begin{definition}[Space Complexity]
  The \textit{space complexity} of a TM $M$ is the function $S(n)$ that is
  the maximum of the space consumption of $M$ on all inputs $\omega$
  of length $n$.
  \begin{center}
    $M$ is $S(n)$ space bounded $\iff\forall$ inputs $\omega$ of length $n$,
    $Space_M\left( \omega \right) \le S(n)$
  \end{center}
\end{definition}

\begin{prop}[Existence of Arbitrary Difficult Problems]
  Suppose a recursive function $f$. We define
  \begin{lstlisting}[mathescape]
  g$\left( x \right)$:
    Simulate $M_x$ on input $x$.
    if $M_x$ does not halt within $f\left(\left|x\right|\right)$ steps, output $0$.
    else output $M_x(x)+1$. // or anything not $M_x(x)$
  endg
  \end{lstlisting}
  $g$ cannot be computed by any $f(n)$ time bound TM.
\end{prop}
\begin{proof}
  Suppose not. Consider the $f(n) $ time bound TM $M_y$ that computes $g$.
  Consider $M_y\left( y \right)$. $M_y$ must halt within $f(\left| y \right|)$
  steps since it is $f(n)$ time bound. Hence $M_y\left( y \right) $ outputs
  $M_y(y)+1$. A contradiction.\qed
\end{proof}

\begin{definition}[Space Constructible Functions]
  A function $S(n)$ is \textit{fully space constructible} if there exists a
  $S(n)$ space bounded TM $M$ such that on all inputs of length $n$, it uses
  space exactly $S(n)$.
\end{definition}

\begin{definition}[Time Constructible Functions]
  A function $T(n)$ is \textit{fully time constructible} if there exists a
  $T(n)$ time bounded TM $M$ such that on all inputs of length $n$, it halts
  and uses time exactly $T(n)$.
\end{definition}

\subsection{Definitions}
\begin{definition}[P]
  A language $L$ is in the class \textit{P} if there exists a
  \textbf{deterministic} TM $M$ and a polynomial time complexity $T(n)$ such
  that $L=L(M)$ and $M$ is $T(n)$ time bounded.
\end{definition}

\begin{definition}[NP]
  A language $L$ is in the class \textit{NP} if there exists a
  \textbf{nondeterministic} TM $M$ and a polynomial time complexity $T(n)$ such
  that $L=L(M)$ and $M$ is $T(n)$ time bounded.
\end{definition}

\begin{definition}[co-NP]
  A language $L$ is in the class \textit{co-NP} if its complement $\bar{L}$ is
  in NP.
\end{definition}

\begin{definition}[Certificate]
  A \textit{certificate} (witness) is a string that certifies the answer to a
  computation, or membership of some string in a language. Intuitively, a
  certificate is the ``answer'' to the problem at hand.
\end{definition}
\begin{prop}
  $L$ is in NP $\iff$ there exists a polynomial $q\left( \cdot \right) $ and a
  polynomial-time bounded TM $M$ such that 
  \[x\in L \iff \exists \text{ a certificate } c, \left| c \right| \le
  q\left(\left| x \right| \right) \text{$M$ accepts the pair $\left( x,c
\right) $}.\]
  That is, we can verify a solution of the problem in polynomial time.
\end{prop}

\subsection{NP-Completeness}
\begin{definition}[NP Hard]
  A language (problem) $L$ is \textit{NP hard} if for every language $L'\in
  NP$ there is a polynomial time reduction of $L'$ to $L$.
\end{definition}

\begin{definition}[NP Complete]
  A language (problem) $L$ is \textit{NP complete} if $L$ is in NP, and it is
  NP hard. That is, for every language $L'\in NP$ there is a polynomial time
  reduction of $L'$ to $L$.
\end{definition}

\begin{definition}[Karp Reducibility]
  A problem $P_1$ is \textit{Karp reducible} to $P_2$, $P_1\le _m^{p} P_2$, if
  there exists a polynomial time computable recursive function $f$ such that
  \[x\in P_1\iff f(x)\in P_2.\]
  We say that $P_1$ is polynomial time, many-one reducible to $P_2$.
\end{definition}

\begin{remark}
  We may further define an alternative notion of reducibility known as a
  \textit{Cook reducibility}.  A problem $P_1$ is \textit{Cook reducible} to
  $P_2$, $P_1\le _T^{p} P_2$, if, given an oracle for $P_2$, there exists a
  polynomial time recursive algorithm that solves $P_1$. This presents another
  type of NP-completeness known as Cook-completeness. In the course and in many
  conventions, we use the more retricted notion of Karp-completeness.
\end{remark}

Since the reducible relation is reflexive and transitive, it is sufficient to
prove that if $L$ is NP complete, $L'\in NP$, $L\le^{p}_m L'$, then $L'$ is NP
complete. We can hence prove the NP-completeness of some problem $P$ with a
known NP-complete problem $P'$ that Karp-reduces to $P$ and that $P$ is NP.

\begin{theorem}
  If  one of the NP complete problems is solvable in polynomial time, all the
  problems in NP are solvable in polynomial time. That is, $P=NP$.
\end{theorem}

\subsection{Known NP-Complete Problems}
\subsubsection{Satisfiability} Suppose a set $U$ of variables and a collection $C$
of clauses over $U$. Does there exist a satisfying truth assignment for $C$?
\begin{example}[Instance of Satisfiability]
  Let $U=\left\{ A,B,C,D,E,F \right\} $ be a set of variables. Consider
  the expression
  \[\left( A\vee B \vee \neg C\vee E\right)\wedge \left( E\vee F\vee \neg A
  \right) \left( C\vee F \right) \wedge \ldots.\]
  Is there a truth assignment to the literals such that the expression is true?
\end{example}
The 3-SAT problem is the Satisfiability problem, with each clause containing
exactly 3 literals.

\subsubsection{3-Dimensional Matching} Suppose $3$ disjoint finite sets $X,Y,Z$ of
equal cardinality $n$,  and a set $S\subseteq X\times Y\times Z $. Does $S$
contain a matching (a subset $S'\subseteq S$ of cardinality $n$ such that for
any two distinct $\left( x_1,y_1,z_1 \right) , \left( x_2,y_2,z_2 \right) \in
S'$, $x_1\neq x_2$, $y_1\neq y_2$, and $z_1\neq z_2$)?\\

\subsubsection{Vertex Cover} Suppose a graph $G=\left( V,E \right) $ and a
positive integer $K\le \left| V \right| $. Is there a vertex cover $V'$ of size
$\le K$?\\

\begin{proof}
  Given a graph $\left( V,E \right) $, guess a $V'\subseteq V$ and verify if
  $\left| V' \right| \le k$ and $\forall \left( v,w \right) \in E$, at least
  one of $v,w$ is in $V'$. Vertex Cover is in NP.\\

  Suppose an instance of 3SAT, with the set of variables $U=\left\{
  x_1,x_2,\ldots,x_{n} \right\} $, the set of clauses $C=\left\{
c_1,c_2,\ldots,c_m \right\} $ where $c_i=\left( l_{i,1}\vee l_{i,2}\vee l_{i,3}
\right) $.

  Consider a vertex cover instance $G=\left( V,E \right) $, where:
  \begin{gather*}
    V=\left\{ u_i,w_i|1\le i\le n \right\} \cup \left\{ z_{j,1},z_{j,2},z_{j,3}
    | 1\le j\le m\right\} \\
    E=\left\{ \left( u_i,w_i \right) | 1\le i\le n \right\} \cup \left\{ \left(
    z_{j,1}, z_{j,2} \right), \left( z_{j,1}, z_{j,3} \right), \left( z_{j,2},
z_{j,3} \right) | 1\le j\le m \right\}\cup \left\{ \left( z_{j,r}, u_i \right)
| l_{j,r}=x_i \right\} \cup \left\{ \left( z_{j,r}, w_i \right) | l_{j,r}=\neg
x_i \right\}\\
  K=2m+n
  \end{gather*}
  Each $u_i$ represents $x_i$, $w_i$ represents $\neq x_i$, and $z_{j,r}$
  represents the $l_{j,r}$. This reduction can be done in polynomial time. We
  want to show that the 3SAT instance has a solution $\iff$ the vertex cover
  instance has a solution.\\

  In a vertex cover, 
  \begin{itemize}
    \item For each $1\le i\le n$, at least one of $u_i, w_i$ is in the cover.
      (A pair representing a variable and its complement)
    \item For each $1\le j\le m$, at least two of $z_{j,1},z_{j,2},z_{j,3}$ is
      in the cover. (A triangle representing each clause)
  \end{itemize}
  Since we have the size bound $K=2m+n$, \textit{exactly} one and two
  (respectively) of the vertices are in the cover.\\

  ($\implies$) Suppose the 3SAT instance is satisfiable. By choosing $u_i$
  $\iff$ $x_i$ is true and $x_i$ $\iff$ $x_i$ is false, and choosing two of
  $z_{j,a}$ such that the one left out represents the true literal, we have a
  vertex cover of $G$.

  ($\impliedby$) Suppose the vertex cover instance has a solution. Consider the
  truth assignment such that $x_i$ is true $\iff$ $u_i$ is in the vertex cover.
  It can be shown that if $z_{j,r}$ is not in the vertex cover then $l_{j,r}$
  must be true.\\

  Hence Vertex Cover is NP Complete.\qed
\end{proof}

\subsubsection{MAX-CUT} Suppose an undirected graph $G=\left( V,E \right) $ and a
positive integer $K\le \left| E \right| $. Does there exist a cut of $G$ with
size $>K$? 

\subsubsection{Clique} Suppose a graph $G=\left( V,E \right) $ and a positive
integer $K\le \left| V \right| $. Does $G$ contain a clique of size $\ge K$ (a
subset larger than $K$ of $V$ that can form a complete graph with edges in
$E$)?

\subsubsection{Independent Set} Suppose a graph $G=\left( V,E \right) $ and a
positive integer $K\le \left| V \right| $. Does $G$ contain an independent set
of size $\ge K$ (a subset larger than $K$ of $V$ such that no distinct pair of
vertices are adjacent)?\\

\begin{proof}
  Given a graph $\left( V,E \right) $, guess a $V'\subseteq V$ and verify if
  $\left| V' \right| \le K$ and that it is a complete graph. Clique is in NP.
  Similarly for Independent Set.\\

  Suppose $G=\left( V,E \right) $. $V'$ is a vertex cover of $G$ $\iff$ $V-V'$
  is an independent set of $G$ $\iff$ $V-V'$ is a clique of $G'$, where $G'$ is
  the complement of $G$. That is, the graph has a vertex cover of size $k$
  $\iff$ $G$ has an independent set of size $\left| V \right| -k$ $\iff$ $G'$
  has a clique of size $\left| V \right| -k$.\\

  Hence Clique/Independent Set is NP Complete.\qed
\end{proof}

\subsubsection{Hamiltonian Circuit} Suppose a graph $G=\left( V,E \right) $. Does
$G$ contain a Hamiltonian circuit (a simple circuit that goes through all
vertices of $G$)?

\subsubsection{Partition} Suppose a finite set $A$ and a size function $s: A\to
\mathbb{R}^{*}_{+}$. Is there a subset $A'\subseteq A$ such that
\[\sum_{a\in A'}s(a)=\sum_{a\in A-A'}s(a)?\]

\subsubsection{Set Cover} Suppose a finite set $A$, a collection of sets from $A$
$\left\{ S_1,S_2,\ldots,S_m \right\} \subseteq \mathcal{P}\left( A \right) $,
and $k\in \mathbb{N}$. Does there exist a subset $Y\subseteq \left\{
1,2,\ldots,m \right\} $, $\left| Y \right| \le k$ such that
\[A\subseteq  \bigcup_{i \in Y} S_i?\]

\subsubsection{Travelling Salesman Problem} Suppose a complete weighted graph
$G=\left( V,E \right) $, and a bound $B\in \mathbb{R}^{*}_{+}$. Is there a
hamiltonian circuit of total weight $\le B$?

\subsection{Modifications to TMs}
\begin{theorem}[Tape Compression]
  Suppose a constant $c\in \mathbb{R}^{*}_{+}$. If a language $L$ is accepted
  by a $S(n)$ space bounded TM $M$ with $k$ tapes, there exists a TM $M'$ with
  $k$ tapes that is $\lceil cS(n)\rceil$ space bounded that accepts $L$.
\end{theorem}
\begin{proof}
  Suppose a TM $M$ that is $S(n)$ space bounded and accepts a language $L$.\\

  Consider the TM $M'$, where each cell of a worktape of $M'$ encodes $m$ cells
  of the corresponding tape of $M$. Then the tape alphabet
  $\Gamma_{M'}={\Gamma_M}^{m}$. Through variable storage in the finite control
  of $M'$, we keep track of which of the $m$ cells is the currently scanned
  cell. Hence $M'$ is able to simulate $M$, and accepts $\iff$ $M$ accepts.\\

  The space used by $M'$ on input $x$ is:
  \[\lceil \frac{Space_M\left( x \right) }{m}\rceil.\]
  Taking $m>\frac{1}{c}$, $M'$ is $\lceil cS(n) \rceil$ space bounded.\qed
\end{proof}

\begin{theorem}[Linear Speedup]
  Suppose a constant $c\in \mathbb{R}^{*}_{+}$. Suppose a language $L$ is accepted
  by a $T(n)$ time bounded TM $M$ with $k\ge 2$ tapes, where
  \[\lim_{n \to \infty} \frac{T(n)}{n}=\infty.\]
  Then there exists a TM $M'$ with $k$ tapes that is $\lceil cT(n)\rceil$ time
  bounded that accepts $L$.
\end{theorem}

\begin{corollary}[Linear Speedup]
  Suppose a constant $c\in \mathbb{R}^{*}_{+}$. Suppose a language $L$ is accepted
  by a $T(n)$ time bounded TM $M$ with $k\ge 2$ tapes, where
  \[T(n)=dn,\qquad d\in \mathbb{R}^{*}_+.\]
  Then there exists a TM $M'$ with $k$ tapes that is $\left( 1+c \right) n$
  time bounded that accepts $L$.
\end{corollary}

\subsection{Complexity Classes}
\begin{notation}[Time and Space Complexity Classes]
  \begin{gather*}
    DSPACE(S(n))=\left\{ L|\exists \text{ DTM }M, L(M)=L\wedge M\text{ is
    }S(n)\text{ space bounded.} \right\} \\
    DTIME(T(n))=\left\{ L|\exists \text{ DTM }M, L(M)=L\wedge M\text{ is
    }T(n)\text{ time bounded.} \right\} \\
    NSPACE(S(n))=\left\{ L|\exists \text{ NTM }M, L(M)=L\wedge M\text{ is
    }S(n)\text{ space bounded.} \right\} \\
    NTIME(T(n))=\left\{ L|\exists \text{ NTM }M, L(M)=L\wedge M\text{ is
    }T(n)\text{ time bounded.} \right\} 
  \end{gather*}
\end{notation}

\begin{theorem}[Relationships between Complexity Classes]
  \begin{gather*}
    DTIME\left( S\left( n \right)  \right) \subseteq DSPACE\left( S\left( n
    \right)  \right) \\
    \text{If }L\in DSPACE\left( S\left( n \right)  \right)\text{ and }S\left(
    n \right) \ge\log{n}\text{, then there exists a constant }c(L)\in
    \mathbb{R_+^{*}}, L\in DTIME\left( c^{S\left( n \right) } \right)\\
    \text{If }L \in NTIME\left( T\left( n \right) \right)\text{, then there
    exists a constant } c(L)\in \mathbb{R}^{*}_+, L\in DTIME\left( c^{T\left( n
    \right) } \right) 
  \end{gather*}
\end{theorem}

\begin{theorem}[Savitch's Theorem]
  If $S(n)\ge \log{n}$, and $S\left( n \right) $ is fully space constructible, then
  \[NSPACE\left( S\left( n \right)  \right) \subseteq DSPACE\left( {\left( S\left( n \right)  \right) }^{2} \right).\] 
\end{theorem}

\begin{theorem}
  The class of languages that can be accepted in polynomial space by a DTM
  $PSPACE$ is equal to the class of languages that can be accepted in
  polynomial space by a NTM $NPSPACE$. $PSPACE=NPSPACE$
\end{theorem}

\subsubsection{Hierachy Theorems}
\begin{theorem}
  If $L$ is accepted by a $S(n)\ge \log{n}$ space bounded TM, then $L$ can be
  accepted by a $S(n)$ space bounded machine which halts on all inputs.
\end{theorem}
Hence, WLOG, we may assume that a TM with a large enough space bound halts on
all inputs.
\begin{theorem}[Space Hierachy]
  Suppose $S_2(n),S_1(n)\ge \log{n}$. If $S_2(n)$ is fully space constructible and
  \[\lim_{n \to \infty} \frac{S_1(n)}{S_2(n)}=0,\] 
  there is a language in $DSPACE\left( S_2\left( n \right)  \right)
  -DSPACE\left( S_1\left( n \right)  \right) $.
\end{theorem}
\begin{proof}
  Consider the TM $M$ which is $S_2\left( n \right) $ space bounded, with a
  fixed number of tapes.
  \begin{lstlisting}[mathescape]
  M$\left(x\right)$:
    if $x$ is of the form $1^k$, reject.
    Mark out space $S_2\left(\left|1^{k}0x\right|\right)$ on the work tape.
    Simulate machine $M_x$ on input $1^{k}0x$.
    if $M_x$ attempts to use more than $S_2\left(\left|1^k 0x\right|\right)$, $M$ rejects.
    if $M_x$ halts and rejects the input, $M$ accepts.
  endM
  \end{lstlisting}
  Let $L=L(M)$. Since $M$ is $S_2(n)$ space bounded, $L\in DSPACE\left(
  S_2\left( n \right)  \right) $.\\

  Suppose by way of contradiction that some TM $M'$ is $S_1(n)$ space bounded
  accepting $L$. There exists another $2$-tape machine $M_x$ accepting $L$
  using space $\le dS_1(n)$ for some constant $d$. WLOG, assume $M_x$ halts on
  all inputs.\\

  Simulation of $M_x$ requires space $cS_1(n)$ for some constant $c$. Hence, let
  $k$ be large enough such that
  \[cS_1\left( \left| 1^{k}0x \right|  \right) <S_2\left( \left| 1^{k}0x
  \right|  \right).\]
  The simulation of $M_x$ must complete, and $M$ accepts $\iff$ $M_x$ did not.
  $M_x$ accepts a different language than $L$. A contradiction. \qed
\end{proof}

\begin{theorem}[Time Hierachy]
  Suppose $T_2(n),T_1(n)\ge \left( 1+\epsilon \right) {n}$. If $T_2(n)$ is
  fully time constructible and
  \[\lim_{n \to \infty} \frac{T_1(n)\log{\left( T_1\left( n \right)  \right) }}{T_2(n)}=0,\] 
  there is a language in $DTIME\left( T_2\left( n \right)  \right) -DTIME\left(
  T_1\left( n \right)  \right) $.
\end{theorem}
\begin{proof}
  Consider a machine $M$ with at least $5$ tapes (or more for fully time
  constructibility of $T_2\left( n \right) $). While running $M$, $M$ keeps
  track of the number of steps taken, and rejects if $T_2(n)$ time is exceeded.
  \begin{lstlisting}[mathescape]
  M$\left(x\right)$:
    if $x$ is of the form $1^k$, reject.
    Simulate $M_x$ on input $1^{k}0x$.
    if $M_x$ halts (does not exceed the time limit) and rejects the input, $M$ accepts.
  \end{lstlisting}
  Let $L=L(M)$. Since $M$ is $O(T_2(n))$ time bounded, $L \in DTIME\left(
  T_2(n) \right) $by linear speedup.

  Suppose by way of contradiction that some TM $M'$ is $T_1(n)$ time bounded
  accepting $L$. There exists another $2$-tape machine $M_x$ accepting $L$
  that is $dT_1(n)\log{T_1(n)}$ time bounded, for some constant $d$.\\

  Simulation of $M_x$ requires time $cT_1(n)\log{T_1(n)}$ for some constant
  $c$. Hence, let $k$ be large enough such that
  \[cT_1\left( \left| 1^{k}0x \right|  \right)\log{\left( T_1\left( \left|
  1^{k}0x \right|  \right)  \right) } <T_2\left( \left| 1^{k}0x \right|
\right).\]
  The simulation of $M_x$ must complete, and $M$ accepts $\iff$ $M_x$ did not.
  $M_x$ accepts a different language than $L$. A contradiction. \qed
\end{proof}
